% Generated draft: requires mathematical and source review.
\documentclass[11pt]{article}
\usepackage[margin=25mm]{geometry}
\usepackage{fontspec}
\setmainfont{Latin Modern Roman}
\usepackage{amsmath,amssymb,mathtools,unicode-math}
\setmathfont{Latin Modern Math}
\usepackage{enumitem,needspace,xurl,hyperref}
\setlength{\parindent}{0pt}
\setlength{\parskip}{5pt}
\setlength{\emergencystretch}{3em}
\newcommand{\N}{\mathbb{N}}
\newcommand{\Z}{\mathbb{Z}}
\newcommand{\Q}{\mathbb{Q}}
\newcommand{\R}{\mathbb{R}}
\newcommand{\C}{\mathbb{C}}
\newcommand{\F}{\mathbb{F}}
\newcommand{\A}{\mathbb{A}}
\newcommand{\PP}{\mathbb{P}}
\newcommand{\E}{\mathbb{E}}
\newcommand{\eps}{\varepsilon}
\newcommand{\dd}{\,\mathrm{d}}
\begin{document}
\section*{UnsolvedMath ID: 10600049}
\textbf{Source problem code:} catalogue entry
\subsection{Definitions and mathematical statement}
Let \(\mathcal L_{\mathrm{virt}}\), \(\mathcal L_{\mathrm{flat}}\) and \(\mathcal L_{\mathrm{free}}\) denote the respective equivalence classes of long virtual, long flat and long free knots, with the standard behavior outside a compact part of the long strand. Virtual diagrams retain classical crossing over/under data and virtual crossings; flat diagrams forget classical over/under information; the free version further forgets the pertinent local crossing data, usually modeled by unlabelled chord diagrams on an oriented line with the free Reidemeister relations. Exact conventions, including orientations and virtualizations, require confirmation from the survey.

The source credits Turaev with a map \(\mathcal L_{\mathrm{flat}}\to\mathcal L_{\mathrm{virt}}\) and asks for a map
\[s:\mathcal L_{\mathrm{free}}\longrightarrow\mathcal L_{\mathrm{virt}}.\]
Its broader request is to construct maps that ``enlarge'' structure. Any proposed diagram rule must descend to the respective equivalence classes. Bare set-theoretic existence imposes no substantive condition: the constant map to the long virtual unknot is already a map. Therefore the intended enrichment property is essential missing data.

For example, if \(F:\mathcal L_{\mathrm{virt}}\to\mathcal L_{\mathrm{free}}\) is the intended forgetful map and the goal is to enhance each free knot without changing its underlying free class, the extra condition would be
\[F\circ s=\operatorname{id}_{\mathcal L_{\mathrm{free}}}.\]
That is a possible precise section problem, and would imply injectivity of \(s\); neither this condition nor injectivity, naturality, nonconstancy or compatibility with long-knot composition is expressly stated in the excerpt. It is not silently substituted for the literal question. Similarly the general enlargement request needs specified source/target categories and compatibility conditions before it becomes a precise construction problem.
\subsection{Short English statement}
Construct a meaningful enrichment from long free knots to long virtual knots, inspired by Turaev's flat-to-virtual map. Since constant maps already exist, the source needs an additional property explaining what counts as enlarging structure.
\subsection{Sources}
\begin{itemize}
\item Local UnsolvedMath registry, permanent ID \(10600049\), exact source code \texttt{catalogue entry}.
\item \emph{Unsolved Problems in Virtual Knot Theory and Combinatorial Knot Theory}, revised survey, problem-list item 49: \url{https://arxiv.org/abs/1409.2823}. Registry access 2026-07-29; source-tex extraction.
\item Turaev is credited with the long-flat to long-virtual construction, but no precise primary citation or properties of that map are supplied.
\end{itemize}
\end{document}
